\documentclass[quick,con]{debate}
\motion{语言的边界是 / 不是人类的边界}
\stance{反方}{不是}
\begin{document}
\debatetitle
\begin{thesisbox}语言记下的是人走过的路。人总是先想到、先做到，词才跟上来；所以人能走到的地方，比语言能说到的地方远。\end{thesisbox}

\section{定义与判准（标准表述）}
\begin{fields}
\field{语言}{有词、有语法、用来表达和交流意思的符号系统。口语、文字、手语都是；音乐、绘画叫语言，只是比喻。}
\field{人类的边界}{人能想到、认识到、做到，并且能传给别人的范围的极限。}
\field{是}{两条边界基本重合：语言到不了的地方，人类在重要的方面也到不了。反例要稳定、成体系、对人重要。}
\field{判准}{语言到不了的地方，人类还能不能稳定地到达。\sub{我方要证明}{在语言之外，人能稳定地思考、做到，而且留得住。}\sub{对方要证明}{人超出本能的那一截以语言为条件；语言之外的，要么到不了，要么留不住。只证明语言有用不够。}}
\field{工具不是边界}{语言是人类最好的工具；好用的工具让人走得快，不划人能走多远。}
\field{对方提偏向定义或判准时的 30 秒口径}{对方把语言扩成一切符号：那题目就成了“人能表达的边界是人能表达的边界”。对方把人类缩成“传下来的东西”：人类是一个个人，一个人到了，人类就到了。}
\end{fields}

\section{我方论点}
\begin{tabularx}{\linewidth}{@{}L{5.2em}Y Y L{6.5em}@{}}\toprule
\thead{标签}&\thead{一句话}&\thead{核心论据}&\thead{出处}\\\midrule
\textbf{想在说前面}&思考不是用语言做的；语言坏了，思想还在；语言没到，人已想到&三名重度语法失语者解括号算式；全面失语者仍能加减、解逻辑题、推断他人想法、欣赏音乐；舍巴林失语后作曲；爱因斯坦：词语在第二阶段才费力去找&Varley 等 2005；Fedorenko \& Varley 2016；Fedorenko 等 2024；Hadamard 1945\\
\textbf{知道的比说的多}&棋感、手上的分寸、医生的直觉，说不清却做得到、教得会&只比划不出声，每下成功率翻一倍；俄语 26 人、英语 24 人辨蓝，准确率 95.7\% 与 96.5\%，俄语者快 124 毫秒&Morgan 等 2015；Winawer 等 2007；波兰尼 1966\\
\bottomrule\end{tabularx}

\section{对方论点 → 一句话反驳}
\begin{tabularx}{\linewidth}{@{}Y Y Y@{}}\toprule
\thead{对方会说}&\thead{我方一句话回}&\thead{对方追问时}\\\midrule
尼加拉瓜第一批手语者答不对错误信念题&证明词帮人学会，不证明人到不了；那门手语是没有语言的孩子造的&“二十五年也替代不了”→ 替代不了的是那套词的教学，而词是人造的\\
石器实验说话翻四倍&比划也翻一倍；没有像样的语言也传了七十多万年&“那是停滞”→ 停住的是速度，不是人\\
轮扁：古之人与其不可传也死矣&“不能以喻”是用话教不会；轮扁七十岁还在斫轮&“儿子学不走”→ 学不走的是用话教的办法\\
失语者都是先有语言的（脚手架）&孩子身上两套系统早就分开；没语言的聋孩子也学会数学&“有延迟”→ 慢不是边界\\
一个人摸到不等于人类走到&人类就是一个个人；词是去接他的，不是先去划他的&“那边界在谁脑子里”→ 在走到那里的人身上\\
5.6：我的语言的界限意味着我的世界的界限&两难：是逻辑，5.641 说哲学的“我”不是人；是经验，失语者就是反例&“他说的是逻辑形式”→ 那就与人类的边界无关\\
\bottomrule\end{tabularx}

\section{必问与必答}
\begin{points}{必问（对方不答就点名、重复、不换题）}
\item 一个中风失语的人，还能不能思考？（全场问题）
\item 尼加拉瓜那门手语，是谁造出来的？
\item 比划算不算您方定义里的语言？
\end{points}
\begin{points}{必答（15 秒正面回答）}
\item 说不出的东西，怎么让下一个人接着往前走？→ 做给他看，手把手；比划也翻一倍，手艺几千年没断。
\item 失语者是不是先有语言的？→ 是；可算的时候不靠语言，孩子身上两套系统早就分开。
\item 说话翻四倍，你们认不认？→ 认，说话更快；快是工具好用，不是边界。
\end{points}

\section{对辩战场概要（本赛制没有自由辩，三场对辩就是交锋主场）}
\begin{tabularx}{\linewidth}{@{}L{6em}Y Y Y@{}}\toprule
\thead{战场}&\thead{开场问题}&\thead{目标}&\thead{收束语}\\\midrule
三辩对辩：思考&失语的人还能不能思考？&语言坏了，思想还在；手语是没语言的人造的&语言缩到零，世界没缩\\
一辩对辩：工具与边界&一个人先到了，这一步算不算人类的？&拆掉“摸到与走到”，钉住“工具不是边界”&词是去接人的，不是先去划人的\\
二辩对辩：稳定&失语者一次次解出算式，算不算稳定？&证明我方例子稳定、成体系、代代有人到&做得到的，比说得出的多\\
\bottomrule\end{tabularx}

\section{如果……就……}
\begin{contingency}
\ifthen{对方把语言定义成一切符号}{指出这是把结论藏进定义，回到有词有语法的中立定义。}
\ifthen{“想在说前面”被打穿}{收缩到“知道的比说的多”，用判准的“稳定到达”解释为什么仍然占优。}
\end{contingency}

分工：一辩守定义判准、为结辩取材，二辩守论点一，三辩守论点二。

\begin{recordcard}
\record{反二质询结束}{质询拿到了什么}{\opt{质询拿到了对方承认失语的人还能思考}\opt{质询里对方回避了失语的人能不能思考}\opt{质询拿到了对方承认比划也教得会}\opt{质询没有拿到明确成果}}{反一立论·开头}
\record{正一立论结束}{对方立论的主干}{\opt{对方立论建立在尼加拉瓜手语者答不对错误信念题上}\opt{对方立论建立在石器实验说话翻四倍上}\opt{对方立论建立在维特根斯坦那句名言上}\opt{对方立论的主干不在以上几条时}}{反二申论·拆正方立论}
\record{正三申论结束}{正三申论主打什么}{\opt{正三申论主打性骚扰这个词}\opt{正三申论主打师傅走了说不出的就断了}\opt{正三申论主打失语者是先有语言的}\opt{正三申论不在以上几条时}}{反三申论·回应正三}
\record{反一质询正三结束}{正方全场最有力的一点}{\opt{正方全场最有力的是尼加拉瓜手语者}\opt{正方全场最有力的是石器实验说话翻四倍}\opt{正方全场最有力的是一个人摸到不等于人类走到}\opt{正方最有力的一点不在以上几条时}}{反三申论·处理正方最强点}
\record{二辩对辩结束}{全场主战场落在}{\opt{全场主战场落在失语者能不能思考}\opt{全场主战场落在说不出的能不能传下去}\opt{主战场不清楚时}}{反一结辩·归纳分歧}
\record{二辩对辩结束}{正方全场最有力的一点}{\opt{正方全场最有力的是尼加拉瓜手语者}\opt{正方全场最有力的是石器实验说话翻四倍}\opt{正方最有力的一点不在以上几条时}}{反一结辩·处理最强点}
\record{二辩对辩结束}{正方全场主打的路线}{\opt{正方全场主打一个人摸到不等于人类走到}\opt{正方全场主打没有词就走不到}\opt{正方全场路线不清楚时}}{反一结辩·封路}
\record{全场}{对方回避了哪些问题}{（写下来）}{申论与结辩里点名}
\end{recordcard}
\end{document}
